\centerline{\bf \Huge Комбинаторная диагностика}

\begin{flushright}
        \scriptsize{Дям дям диди диди}
\end{flushright}

\problem{1}
В математическом соревновании участвовали $200$ школьников, которым было предложено $6$ задач. Каждую задачу правильно решили не менее $120$ участников. Докажите, что найдутся два участника, которые в совокупности решили все шесть задач.
% Источник: IMC 2002; Yufei Zhao, Counting in Two Ways, Problem 2.
% https://yufeizhao.com/olympiad/doublecounting_mop.pdf

\problem{2}
В английском клубе вечером собрались $n$ его членов, где $n\geqslant 3$. Каждый принёс сок того вида, который он предпочитает, в том количестве, которое собирается выпить. В любой момент любые три члена клуба могут сесть за один столик и выпить некоторое количество сока, причём все трое обязательно выпивают поровну каждый своего сока.

Докажите, что все члены клуба смогут полностью выпить принесённый сок тогда и только тогда, когда количество сока, принесённое каждым из них, не превосходит одной трети общего количества сока.
% Источник: Московская математическая олимпиада, 2016, 11 класс, задача 4;
% И.\,В. Яковлев, листок «Принцип крайнего», задача 7.
% https://mathus.ru/math/prinkrai.pdf

\problem{3}
В стране некоторые пары городов соединены односторонними прямыми авиарейсами; между любыми двумя городами имеется не более одного рейса. Будем говорить, что город $A$ доступен из города $B$, если из $B$ можно долететь в $A$, возможно, с пересадками. Считается также, что каждый город доступен из самого себя.

Известно, что для любых двух городов $P$ и $Q$ существует город $R$, доступный и из $P$, и из $Q$. Докажите, что существует город, доступный из всех городов страны.
% Источник: заключительный этап ВсОШ, 2017, 9 класс, задача 1;
% И.\,В. Яковлев, листок «Графы», задача 13.
% https://mathus.ru/math/graphs.pdf

\problem{4}
Квадрат со стороной $1$ двумя различными способами разбит на $n$ многоугольников равной площади. Докажите, что в квадрате можно выбрать $n$ точек так, чтобы в каждом многоугольнике каждого из двух разбиений находилось ровно по одной выбранной точке.
% Источник: СУНЦ МГУ, листок «Теорема Холла», 13 ноября 2013 года, задача 2.
% https://math.mosolymp.ru/upload/files/2014/AESC/rounding/2013-11-13-hall-theorem.pdf

\problem{5}
Таблица $10\times 10$ заполнена по правилам игры «Сапёр»: в некоторых клетках стоят мины, а в каждой из остальных клеток записано количество мин в клетках, соседних с ней по стороне или по вершине.

Все первоначальные мины убрали, а во все клетки, в которых мин не было, поставили мины. После этого числа в свободных клетках записали заново по тем же правилам. Может ли сумма всех чисел в таблице увеличиться?
% Источник: Турнир городов, 2013, старшая группа, задача 1;
% И.\,В. Яковлев, листок «Инварианты», задача 20.
% https://mathus.ru/math/invariant.pdf